% ── sections/05_conclusions.tex ───────────────────────────────────────────────

\section{Conclusion}

This work presented a systematic investigation of cancer driver gene prediction
using multilayer graph neural networks across six PPI databases and 64-dimensional
multi-omics features. We faithfully reproduced the EMGNN benchmark (mean
AUPR~=~0.743, 18 runs on CPDB), identified that BatchNorm is harmful in full-batch
graph settings ($-$4.2\% AUPR), and demonstrated that multi-network training with
learnable importance weights achieves AUPR~=~0.8067 (+5.9\%), with gain
decomposition revealing data integration contributes +5.4\% while architectural
modifications contribute only +1.0\%. The six-network baseline is robust across
five seeds (0.8019~$\pm$~0.0050).

Our core finding is that heterophily-aware gating provides the largest single
improvement among nine evaluated techniques (AUPR~=~\textbf{0.824~$\pm$~0.004},
+2.8\%, $p<0.01$). Cross-network attention (+1.5\%) and DropEdge (+0.6\%) provide
moderate, consistent gains. Two important negative results emerged: Focal Loss
degrades performance ($-$5.1\%), demonstrating that label smoothing alone is
sufficient for calibration, and GraphMAE pretraining provides no benefit
($-$0.2\%), indicating that 64-dimensional multi-omics features are already
sufficiently informative. Integrated Gradients attribution and gene set enrichment
validate predictions against 28 significant Hallmark pathways, confirming that the
model has learned biologically meaningful patterns.

These findings establish that structure-aware innovations (heterophily-aware
gating, cross-network attention) outperform loss-function modifications (Focal
Loss) and pretraining (GraphMAE) for this task, and that multi-network data
integration---not architectural complexity---drives the majority of predictive
gains. Source code and experiment scripts are publicly available at
\url{https://github.com/EthyleneC2H4/NUS-Capstone}.
